\documentclass[a4paper]{article}
\usepackage[utf8x]{inputenc}    
\usepackage[T1]{fontenc}
\usepackage[francais,bloc,nowatermark]{automultiplechoice}    
% [nowatermark] pour enlever le texte "PROJET" et le pied de page
\usepackage{multicol}
\usepackage{eurosym}
\usepackage{fancyhdr,amssymb,amsmath}

%===================
\begin{document}
%===================

%%%%%%%%%%%%%%%%%%%%%%%%%
%% début des questions %%%%
%%%%%%%%%%%%%%%%%%%%%%%%%

%%% On répartie les quesitons dans différents groupes.
% ---------------------------------
%%% Questions groupe : "catégorie1"
%----------------------------------
\element{categorie1}{
  \begin{question}{calcul littéral 1}    
% Question ouverte   

Une expérience aléatoire consiste à lancer un dé à six faces et regarder le résultat obtenu. L'univers associé à cette expérience est l'ensemble de toutes les issues possibles : Ici $\Omega=\{1 ; 2 ; 3 ; 4;5 ; 6\}$

On considère le jeu suivant :
\begin{itemize}
\item Si le résultat est pair, on gagne $2$ \euro;
\item Si le résultat est 1 , on gagne $3$ \euro;
\item Si le résultat est 3 ou 5 , on perd $6$ \euro.
\end{itemize}
On note $X$ la variable aléatoire représentant le gain ou la perte en euro du joueur.

\begin{enumerate}
\item Quelle sont les valeurs possibles pour $X$ ? On notera les possibles valeurs $x_1; x_2,\cdots$ etc.
\item Donner la loi de probabilité de $X$ à l'aide d'un tableau.
\item Déterminer l'espérance de $X$, notée $E(X)$.
\item Le jeu est-il équitable? 


\end{enumerate}
   \AMCOpen{lines=7,dots=false}{\wrongchoice{NA}\scoring{0}\wrongchoice{AB}\scoring{1}\wrongchoice{B}\scoring{2}\correctchoice{TB}\scoring{3}}
% AB rapporte 1 point
% B rapporte 2 points
% TB rapporte 3 points     
  \end{question}
}

%
\element{categorie2}{
  \begin{question}{calcul littéral 2}  
% Question ouverte  
On note $X$ la variable aléatoire qui, à chaque jour, associe le nombre de véhicules neufs vendus par un concessionnaire. Sa loi de probabilité est donnée par le tableau suivant :

\begin{center}
\begin{tabular}{|c|c|c|c|c|}
\hline$x_{i}$ & 0 & 1 & 2 & 3 \\
\hline$\mathbb{P}\left(X=x_{i}\right)$ & 0,45 & 0,3 & 0,15 & $\ldots \ldots$ \\
\hline
\end{tabular}
\end{center}

\begin{enumerate}
\item Donner la probabilité $\mathbb{P}(X=1)$. Interpréter ce résultat.
\item Quelle est la probabilité que le concessionnaire vende trois véhicules dans la journée?
\item Calculer $\mathbb{P}(X \geqslant 1)$. Interpréter ce résultat.
\item Combien de véhicules par jour vend en moyenne le concessionnaire?
\end{enumerate}


   \AMCOpen{lines=11,dots=false}{\wrongchoice{NA}\scoring{0}\wrongchoice{AB}\scoring{1}\wrongchoice{B}\scoring{2}\correctchoice{TB}\scoring{3}}
% AB rapporte 1 point
% B rapporte 2 points
% TB rapporte 3 points         
  \end{question}
}


%%%%%%%%%%%%%%%%%%%%%%%%%
%% fin des questions %%%%
%%%%%%%%%%%%%%%%%%%%%%%%%


%%%%%%%%%%%%%%%%%%%%%%%%%%%%
%%% fabrication des copies%%
%%%%%%%%%%%%%%%%%%%%%%%%%%%%
\exemplaire{18}{    

%%% debut de l'entête des copies :    
	%%% première ligne

	\begin{minipage}{1\linewidth}
	  \centering\large\bf DS TSTMG sur les variables aléatoires discrètes 
	\end{minipage}

	%%% codage des numéros d'étudiants sur 2 chiffres par les candidats et encadré d'écriture manuelle des noms.
	\vspace*{3mm}
	{\setlength{\parindent}{0pt}\hspace*{\fill}
	\begin{minipage}[b]{12cm}
	 

	 \hfill\champnom{\fbox{
	 \begin{minipage}{1\linewidth}
		\vspace*{3mm}
		NOM - Prénom - Classe :
		\vspace*{3mm}
	 \end{minipage}
	    }
	   }
	\hfill\vspace{1ex}
 %texte de présentation 
\begin{center}\em

Durée : 55 minutes.

  Aucun document n'est autorisé. La calculatrice est autorisée.
%  Les questions faisant apparaître le symbole \multiSymbole{} peuvent
%  présenter une ou plusieurs bonnes réponses. Les autres ont une unique bonne réponse.
%  Les réponses fausses ou incohérentes retirent des points.

\end{center}

	\end{minipage}\hspace*{\fill}
	}

\vspace{4ex}

%%% Fin de l'entête
%%%%%%%%%%%%%%%%%%%%%%%

%% Composition des QCM
%%% On mélange des questions par groupe
% On efface le {tout} précédent
\cleargroup{tout}

% On mélange par "catégorie", et on prélève aléatoirement [n] questions 
% du lot {categorie-k} vers la copie complète {tout}
% puis on restitue la copie.
% Si on met [0], toutes les questions sont prises.
\melangegroupe{categorie1}\copygroup[0]{categorie1}{tout}
\melangegroupe{categorie2}\copygroup[0]{categorie2}{tout}
%\melangegroupe{categorie3}\copygroup[0]{categorie3}{tout}
%\melangegroupe{categorie4}\copygroup[0]{categorie4}{tout}

\restituegroupe{tout}

% Saut d'une page si le sujet comporte un nombre impair de pages.
% Ainsi il est possible d'imprimer les sujets en recto-verso.

\AMCcleardoublepage    

%%% Fin de la création des questionnaires

}
% Fin de exemplaires{k}{ 

\end{document}

